\documentclass[10pt,a4paper,fontset=fandol]{ctexart}

\usepackage[
  margin=1.8cm,
  top=1.55cm,
  bottom=1.55cm
]{geometry}

\usepackage{
  fontspec,
  xcolor,
  enumitem,
  titlesec,
  hyperref,
  parskip
}

% English font
\setmainfont{TeX Gyre Termes}

% Chinese fonts are automatically
% configured by ctex (fontset=fandol).

\definecolor{navy}{HTML}{183650}

\hypersetup{
  colorlinks=true,
  urlcolor=navy
}

\setlength{\parindent}{0pt}
\setlength{\parskip}{3pt}

\setlist[itemize]{
  leftmargin=1.55em,
  topsep=2pt,
  itemsep=1pt,
  parsep=0pt
}

\titleformat{\section}
  {\large\bfseries\color{navy}}
  {}
  {0pt}
  {}
  [\vspace{-4pt}\color{navy}\titlerule]

\titlespacing*{\section}
  {0pt}{10pt}{6pt}
\newcommand{\entry}[3]{%
  \textbf{#1}\hfill{\small #2}\\%
  {\emph{#3}}\par%
}

\begin{document}
\begin{center}
{\LARGE\bfseries 舒屹扬}\\[4pt]
{\normalsize 数学方向本科生｜学术简历}\\[2pt]
中国科学技术大学 · 少年班学院\\[2pt]
中国·安徽合肥 \quad\textbar\quad \href{mailto:shuyiy@mail.ustc.edu.cn}{shuyiy@mail.ustc.edu.cn} \quad\textbar\quad \href{https://nicholasshu.com}{个人主页}
\end{center}
\vspace{-7pt}
\section*{研究兴趣}
微分几何、几何分析、凸几何与 Kakeya 型问题。主要关注凸 Kakeya 集及相关变形问题、极值几何构型，以及凸性和几何约束与面积、直径、周长等几何量之间的关系。
\section*{教育背景}
\entry{中国科学技术大学}{2025.09--至今}{少年班学院 · 数学方向本科生（纯数学方向）}
2025 年获得中国科学技术大学少年班及创新试点班招生选拔 A 档资格并入学。
\entry{安徽省安庆市第一中学}{2023--2025}{安徽安庆}
2023、2024 年获高中数学联赛安徽省二等奖；2023 年通过西安交通大学少年班选拔，数学单科成绩全国前 20。
\section*{科研经历}
\entry{凸 Kakeya 集的研究}{2026.09--至今}{指导教师：王作勤教授，中国科学技术大学}
\begin{itemize}
\item 研究在每个方向上均包含单位线段的凸集，关注其面积、直径、周长等几何量之间的定量关系。
\item 探索极值构型、相关几何不等式及凸体在极限过程中的行为。
\item 学习并应用支撑超平面、支撑函数、凸包等凸几何工具。
\end{itemize}
\section*{独立数学论文写作}
\textbf{《双曲平面的一些几何性质》}\hfill{\small 2026}\\
\emph{独立撰写的数学手稿；发表在校数学院的学生刊物《蛙鸣》第68期。}
\begin{itemize}
\item 研究了双曲几何中 Poincaré 上半平面模型的性质，包括其与 Hilbert 公理系统的相容性，以及该模型上的变换群。
\item 探讨了双曲三角形的性质，包括双曲正弦定理和余弦定理，并研究了 Desargues 定理和 Pappus 定理。
\end{itemize}
\section*{荣誉与奖项}
\begin{itemize}
\item \textbf{2026 年：} 中国科学技术大学 Robogame 比赛第一名。
\item \textbf{2025 年：} CMC 数学竞赛省二等奖。
\item \textbf{2025 年：} 中国科学技术大学少年班及创新试点班招生选拔 A 档资格。
\item \textbf{2023、2024 年：} 高中数学联赛安徽省二等奖（各一次）。
\item \textbf{2023 年：} 西安交通大学少年班选拔考试数学单科全国前 20。
\item \textbf{2023 年：} 丘成桐中学生数学夏令营顺利结业。
\item \textbf{2019、2020 年：} 信息学竞赛省二等奖（各一次）。
\item 高中期间曾获优秀学生干部、优秀团干部等荣誉。
\item \textbf{2025 年：} 中国科学技术大学九章杯三等奖。
\end{itemize}
\section*{数学学习背景}
\textbf{基础课程及领域：} 数学分析、线性代数、常微分方程、经典微分几何、凸分析。\\
\textbf{高阶阅读与自主学习：} 光滑流形、微分拓扑、几何分析、偏微分方程、极小曲面、微分形式、de Rham 上同调。\\
\emph{说明：以上兼含正式课程与自主学习内容，具体课程成绩尚未填写。}
\section*{学生工作与技能}
\textbf{学生工作：} 少年班学院班级团支部书记（2025--至今）；少年班学院学生会学术实践部干事（2026--至今）。\\
\textbf{技术工具：} LaTeX、XeLaTeX、Beamer、SolidWorks。\quad\textbf{语言：} 中文（母语）、英语（水平待补充）。
\end{document}
